% © 2025 Ing. David Jaroš — CC BY-NC-ND 4.0
%
% This work is licensed under a Creative Commons Attribution-NonCommercial-NoDerivatives
% 4.0 International License (CC BY-NC-ND 4.0).

\documentclass[11pt,a4paper]{article}
\usepackage{amsmath, amssymb, amsthm}
\usepackage{hyperref}
\usepackage{geometry}
\geometry{margin=2.5cm}

\newtheorem{theorem}{Theorem}
\newtheorem{proposition}{Proposition}
\newtheorem{definition}{Definition}
\newtheorem{remark}{Remark}

\title{Canonical Action Principle\\
\large Unified Biquaternion Theory — Canonical Theory Series}
\author{Ing.\ David Jaroš}
\date{2025}

\begin{document}
\maketitle

\begin{abstract}
This document presents the canonical form of the UBT action principle.
It does not contain new derivations; it collects and cross-references
the action-principle statements already present in the repository.
The UBT action generalizes the Einstein--Hilbert action with Standard Model
matter content by embedding all fields in a single biquaternionic field $\Theta$.
\end{abstract}

\section{The UBT Action}

The action of Unified Biquaternion Theory is:
\begin{equation}\label{eq:S}
  S[\Theta]
  \;=\;
  \int_{\mathcal{M}} \mathrm{d}^4x\,\sqrt{-g}\,\mathcal{L}_{\mathrm{UBT}}[\Theta,\partial\Theta],
\end{equation}
where $\mathcal{M}$ is the real spacetime manifold, $g = \det(g_{\mu\nu})$,
and the Lagrangian density $\mathcal{L}_{\mathrm{UBT}}$ decomposes as:
\begin{equation}\label{eq:Ltotal}
  \mathcal{L}_{\mathrm{UBT}}
  \;=\;
  \underbrace{\frac{1}{2\kappa}\,R}_{\mathcal{L}_{\mathrm{grav}}}
  \;+\;
  \underbrace{\mathrm{Tr}\bigl[(D_\mu\Theta)^\dagger(D^\mu\Theta)\bigr]}_{\mathcal{L}_{\mathrm{kin}}}
  \;-\;
  \underbrace{\frac{1}{4}\,F_{\mu\nu}^a F^{a\,\mu\nu}}_{\mathcal{L}_{\mathrm{gauge}}}
  \;-\;
  \underbrace{V(\Theta)}_{\mathcal{L}_{\mathrm{pot}}},
\end{equation}
with:
\begin{itemize}
  \item $R = R[g]$ — Ricci scalar derived from the projected metric
        $g_{\mu\nu} = \mathrm{Re}(\mathcal{G}_{\mu\nu}[\Theta])$
  \item $\kappa = 8\pi G/c^4$ — gravitational coupling constant
  \item $D_\mu = \partial_\mu + \Gamma_\mu + igA_\mu^a T_a$ — full covariant derivative
        (Levi-Civita $\Gamma$ plus gauge connection $A$)
  \item $F_{\mu\nu}^a$ — field-strength tensor for each gauge group generator $T_a$
  \item $V(\Theta)$ — potential encoding mass terms and symmetry breaking
\end{itemize}

\noindent\textbf{Primary source}:
\texttt{unified\_biquaternion\_theory/ubt\_appendix\_1\_biquaternion\_gravity.tex}, \S1;
\texttt{canonical/interactions/qed.tex}, eq.\ (1);
\texttt{canonical/geometry/gr\_as\_limit.tex}, \S3.

\section{Variational Principle and Field Equations}

The equations of motion follow from requiring stationarity:
\begin{equation}
  \delta S[\Theta] = 0
  \quad \text{for all compactly supported variations } \delta\Theta.
\end{equation}

\subsection{Gravitational Sector}

Varying \eqref{eq:S} with respect to the metric $g^{\mu\nu}$ (equivalently,
with respect to the real part of $\Theta$ entering $\mathcal{G}_{\mu\nu}$)
yields the Einstein field equations:
\begin{equation}
  G_{\mu\nu} \;=\; R_{\mu\nu} - \tfrac{1}{2}g_{\mu\nu}R \;=\; \kappa\,T_{\mu\nu},
\end{equation}
where $T_{\mu\nu}$ is the stress-energy tensor derived from the matter
and gauge sectors. \textbf{Status: PROVED [L1].}

\noindent\textbf{Source}: \texttt{canonical/geometry/gr\_as\_limit.tex}, \S3--\S4;
\texttt{canonical/geometry/stress\_energy.tex};
\texttt{consolidation\_project/GR\_closure/GR\_chain\_summary.tex}, Step 5.

\subsection{Matter/Gauge Sector}

Varying with respect to $\Theta^\dagger$ yields the matter field equations:
\begin{equation}
  D_\mu D^\mu\,\Theta \;=\; \frac{\partial V}{\partial \Theta^\dagger},
\end{equation}
which reduces to the Dirac equation in the appropriate limit.

Varying with respect to $A_\mu^a$ yields the Yang--Mills equations:
\begin{equation}
  D_\mu F^{a\,\mu\nu} \;=\; j^{a\,\nu},
\end{equation}
with gauge current $j^{a\,\nu}$.
In the U(1) abelian limit this becomes $\partial_\mu F^{\mu\nu} = j^\nu$ —
the inhomogeneous Maxwell equations.
\textbf{Status: Maxwell sector PROVED [L1]. Full Yang--Mills sector PARTIAL [L2].}

\noindent\textbf{Source}: \texttt{canonical/interactions/qed.tex}, \S3;
\texttt{canonical/bridges/Maxwell\_limit\_bridge.tex};
\texttt{consolidation\_project/appendix\_C\_electromagnetism\_gauge\_consolidated.tex}.

\section{Biquaternionic Form of the Action}

Before projection to real fields, the action can be written in fully
biquaternionic form:
\begin{equation}
  S[\Theta]
  \;=\;
  \int \mathrm{d}^4x\,\mathrm{d}\psi\;\mathrm{Re}\,\mathrm{Tr}\Bigl[
    \bigl(\nabla^\dagger\Theta\bigr)^*\bigl(\nabla^\dagger\Theta\bigr)
    \;-\;
    V(\Theta)
  \Bigr],
\end{equation}
where integration over $\psi$ (imaginary time) is over the compact fiber
$[0, 2\pi R_\psi]$ and $\nabla^\dagger$ is the biquaternionic conjugate gradient.
This form makes the master equation
$\nabla^\dagger\nabla\,\Theta = \kappa\,\mathcal{T}$ manifest.

\noindent\textbf{Source}: \texttt{canonical/fields/theta\_field.tex}, \S3.

\section{Reduction to Known Theories}

\begin{center}
\begin{tabular}{lll}
\hline
\textbf{Limit} & \textbf{Result} & \textbf{Status} \\
\hline
$\psi \to 0$, real projection & Einstein--Hilbert action + SM & PROVED [L1] \\
U(1) gauge sector, flat spacetime & QED Lagrangian & PROVED [L1] \\
U(1) + constant phase, flat & Maxwell Lagrangian & PROVED [L1] \\
SU(3)$\times$SU(2)$_L\times$U(1)$_Y$ & Standard Model action & PROVED [L1] \\
\hline
\end{tabular}
\end{center}

\noindent For detailed derivation chains see:
\texttt{canonical/bridges/GR\_chain\_bridge.tex},
\texttt{canonical/bridges/QED\_limit\_bridge.tex},
\texttt{canonical/bridges/Maxwell\_limit\_bridge.tex},
\texttt{canonical/bridges/gauge\_emergence\_bridge.tex}.

\section{Cross-References}

\begin{itemize}
  \item Axioms: \texttt{THEORY/canonical/canonical\_axioms.tex}
  \item Metric: \texttt{THEORY/canonical/canonical\_metric\_definition.tex}
  \item Field equations: \texttt{THEORY/canonical/canonical\_field\_equations.tex}
  \item Projection rules: \texttt{THEORY/canonical/canonical\_projection\_rules.tex}
\end{itemize}

\end{document}
